% GENERATED FILE. Edit canonical lessons, then run npm run generate:books.
% canonical-source-hash: fnv1a64:5f431729007a8481
% canonical-lessons: MR-C211-apayas, MR-C211-anubhav, MR-C211-nivad, MR-C211-nirnay, MR-C211-sakyata

\chapter{Failure, Experience, Choice, Decision, Possibility}
\label{ch:mr-failure-experience-choice-decision-possibility}
\hlchaptermodality{211}

\begin{chapteropening}
\textbf{By the end of this chapter:} \emph{I can say failure, experience, a choice, a decision and a possibility.}

Complete the last lesson of chapter 211: I can say failure, experience, a choice, a decision and a possibility.
\end{chapteropening}

\section[\texorpdfstring{\mr{अपयश} (apayaś)}{अपयश (apayaś)}]{\mr{अपयश} (apayaś) — failure}

\label{lesson:MR-C211-apayas}



\begin{quote}
Before the new one: say the Marathi for a worker, then the Marathi for an employee.
\end{quote}

\subsection*{You'll want to know: \mr{अपयश}}
\textbf{\mr{अपयश}} — \emph{apayaś} — \textquotedblleft{}failure\textquotedblright{}.

\begin{grammarlens}[title={What you've built}]
One more word for this chapter.
\end{grammarlens}

\subsection*{Your turn}
\begin{itemize}
\raggedright
  \item \emph{Say it:} \emph{apayaś}
  \item \emph{Say it:} \emph{apayaś}, once more
  \item \emph{Say it:} say \emph{apayaś}
  \item \emph{Recall:} say the Marathi for a worker, then the Marathi for an employee, then say \emph{apayaś} again
\end{itemize}

\begin{tcolorbox}[breakable,colback=teal!4,colframe=teal!35!black,title={Before you move on}]
What does \mr{अपयश} mean? (\textquotedblleft{}Failure\textquotedblright{}.) Say it once more.
\end{tcolorbox}
\section[\texorpdfstring{\mr{अनुभव} (anubhav)}{अनुभव (anubhav)}]{\mr{अनुभव} (anubhav) — experience}

\label{lesson:MR-C211-anubhav}



\begin{quote}
Before the new one: say the Marathi for an employee, then the Marathi for failure.
\end{quote}

\subsection*{You'll want to know: \mr{अनुभव}}
\textbf{\mr{अनुभव}} — \emph{anubhav} — \textquotedblleft{}experience\textquotedblright{}.

\begin{grammarlens}[title={What you've built}]
One more word for this chapter.
\end{grammarlens}

\subsection*{Your turn}
\begin{itemize}
\raggedright
  \item \emph{Say it:} \emph{anubhav}
  \item \emph{Say it:} \emph{anubhav}, once more
  \item \emph{Say it:} say \emph{anubhav}
  \item \emph{Recall:} say the Marathi for an employee, then the Marathi for failure, then say \emph{anubhav} again
\end{itemize}

\begin{tcolorbox}[breakable,colback=teal!4,colframe=teal!35!black,title={Before you move on}]
What does \mr{अनुभव} mean? (\textquotedblleft{}Experience\textquotedblright{}.) Say it once more.
\end{tcolorbox}
\section[\texorpdfstring{\mr{निवड} (nivaḍ)}{निवड (nivaḍ)}]{\mr{निवड} (nivaḍ) — a choice}

\label{lesson:MR-C211-nivad}



\begin{quote}
Before the new one: say the Marathi for failure, then the Marathi for experience.
\end{quote}

\subsection*{You'll want to know: \mr{निवड}}
\textbf{\mr{निवड}} — \emph{nivaḍ} — \textquotedblleft{}a choice\textquotedblright{}.

\begin{grammarlens}[title={What you've built}]
One more word for this chapter.
\end{grammarlens}

\subsection*{Your turn}
\begin{itemize}
\raggedright
  \item \emph{Say it:} \emph{nivaḍ}
  \item \emph{Say it:} \emph{nivaḍ}, once more
  \item \emph{Say it:} say \emph{nivaḍ}
  \item \emph{Recall:} say the Marathi for failure, then the Marathi for experience, then say \emph{nivaḍ} again
\end{itemize}

\begin{tcolorbox}[breakable,colback=teal!4,colframe=teal!35!black,title={Before you move on}]
What does \mr{निवड} mean? (\textquotedblleft{}A choice\textquotedblright{}.) Say it once more.
\end{tcolorbox}
\section[\texorpdfstring{\mr{निर्णय} (nirṇay)}{निर्णय (nirṇay)}]{\mr{निर्णय} (nirṇay) — a decision}

\label{lesson:MR-C211-nirnay}



\begin{quote}
Before the new one: say the Marathi for experience, then the Marathi for a choice.
\end{quote}

\subsection*{You'll want to know: \mr{निर्णय}}
\textbf{\mr{निर्णय}} — \emph{nirṇay} — \textquotedblleft{}a decision\textquotedblright{}.

\begin{grammarlens}[title={What you've built}]
One more word for this chapter.
\end{grammarlens}

\subsection*{Your turn}
\begin{itemize}
\raggedright
  \item \emph{Say it:} \emph{nirṇay}
  \item \emph{Say it:} \emph{nirṇay}, once more
  \item \emph{Say it:} say \emph{nirṇay}
  \item \emph{Recall:} say the Marathi for experience, then the Marathi for a choice, then say \emph{nirṇay} again
\end{itemize}

\begin{tcolorbox}[breakable,colback=teal!4,colframe=teal!35!black,title={Before you move on}]
What does \mr{निर्णय} mean? (\textquotedblleft{}A decision\textquotedblright{}.) Say it once more.
\end{tcolorbox}
\section[\texorpdfstring{\mr{शक्यता} (śakyatā)}{शक्यता (śakyatā)}]{\mr{शक्यता} (śakyatā) — a possibility}

\label{lesson:MR-C211-sakyata}



\begin{quote}
Before the new one: say the Marathi for a choice, then the Marathi for a decision.
\end{quote}

\subsection*{You'll want to know: \mr{शक्यता}}
\textbf{\mr{शक्यता}} — \emph{śakyatā} — \textquotedblleft{}a possibility\textquotedblright{}.

\begin{grammarlens}[title={What you've built}]
One more word for this chapter.
\end{grammarlens}

\subsection*{Your turn}
\begin{itemize}
\raggedright
  \item \emph{Say it:} \emph{śakyatā}
  \item \emph{Say it:} \emph{śakyatā}, once more
  \item \emph{Say it:} say \emph{śakyatā}
  \item \emph{Recall:} say the Marathi for a choice, then the Marathi for a decision, then say \emph{śakyatā} again
\end{itemize}

\begin{tcolorbox}[breakable,colback=teal!4,colframe=teal!35!black,title={Before you move on}]
What does \mr{शक्यता} mean? (\textquotedblleft{}A possibility\textquotedblright{}.) Say it once more.
\end{tcolorbox}
